\section{Version A: implementation when capital inherits the intensity of the final good}\label{sec:mkt:A}

This section sets the instruments of Table~\ref{tab:mkt:instruments} under Version~A,
$\Phi^I=\Omega\phi^I$, and shows that the competitive equilibrium of
Section~\ref{sec:mkt:setup:problems} then reproduces the allocation of Problem~\ref{prob:sp:A}.
Section~\ref{sec:mkt:B} does the same under Version~B. The setup of Section~\ref{sec:mkt:setup} is
presupposed throughout; nothing in Section~\ref{sec:sp:B} is.


\subsection{What has to be corrected}\label{sec:mkt:A:what}

Under Version~A the accounting closes through
\begin{align}
\Phi^I &= \Omega\phi^I,
&
\Omega &= \frac{R}{\mathcal D+\phi^IG}\;\ge\;0
\quad\text{automatically},
&
\mathcal D &= \omega^YY+\omega^{N}C^N+\omega^{D}C^D+\omega^WC^W+\omega^CC,
\label{eq:mkt:A:closure}
\end{align}
which is~\eqref{eq:sp:A:Omega} and~\eqref{eq:sp:setup:Omega-phi}. Recall from
Section~\ref{sec:sp:A:foc} the two objects the planner's conditions are written in terms of,
\begin{align}
\chi &\equiv \frac{\Phi^IG}{R}\in\left[0,1\right),
&
\zeta &= \chi\left(p^W-p^M\right),
&
\Delta &= \zeta\Omega,
\label{eq:mkt:A:zeta}
\end{align}
the share of material input embodied in new capital, the multiplier on the intensity
definition~\eqref{eq:sp:A:intensity}, and the composition wedge it generates.

Three margins are unpriced in the equilibrium of Definition~\ref{def:mkt:equilibrium} when the
instruments are zero, and each corresponds to a block of Table~\ref{tab:mkt:instruments}.

\smalltitle{(i) Mass drawn into production.} By the ledger, a tonne entering production leaves as waste
exactly once, and collecting, treating and releasing it is costly. The final-good firm pays $P^R$ for
that tonne and nothing else; the mine pays nothing for the overburden its extraction displaces. The
correction is a charge at the point of entry.

\smalltitle{(ii) Emissions.} The waste-management sector releases $\left(1-\varpi+d^W\varpi\right)W$
tonnes to the environment and recyclers withdraw $d^WR^R$ of it, and neither is priced. The correction
is an emission tax with a credit for recovery.

\smalltitle{(iii) The composition of final-good spending.} This is the margin specific to Version~A.
Spending one more unit of the final good on activity $j$ raises $\mathcal D$ by $\omega^j$, dilutes
$\Omega$ by~\eqref{eq:mkt:A:closure}, lowers $\Phi^I$, and so diverts mass out of the durable stock and
into today's waste flow. Total mass is unaffected; its \emph{timing} is not. No private agent sees this
effect, since $\Omega$ appears in nobody's problem, and the social value of the shift is $\Delta$ per
unit of the base, as~\eqref{eq:sp:A:foc:Omega} established.


\subsection{The instruments}\label{sec:mkt:A:instruments}

\begin{proposition}[Implementation, Version~A]\label{prop:mkt:A:implementation}
Let $\left\{C,I,N,D,\varpi,K^R,W,\Omega\right\}$ solve Problem~\ref{prob:sp:A} with associated shadow
prices $\left\{\lambda,q,p^S,p^X,p^P,p^M,p^W,\zeta\right\}$, and set
\begin{align}
\tau^R &= p^W-\zeta,
&
\tau^{N,S} &= \Omega^{N,S}p^W,
&
\tau^{D,S} &= \Omega^{D,S}p^W,
\label{eq:mkt:A:mass}
\\
\tau^E &= p^P,
&
s^R &= d^Wp^P,
&
\tau^j &= \Delta\omega^j,\quad j\in\left\{Y,C,N,D,W\right\},
\label{eq:mkt:A:emission}
\\
\tau^K &= -\Phi^I\left(1-\chi\right)\left(p^W-p^M\right),
&&&&
\label{eq:mkt:A:tauK}
\end{align}
with $s^W$ from the zero-profit condition~\eqref{eq:mkt:setup:zeroprofit} and $T$
from~\eqref{eq:mkt:setup:govbudget}. Then the competitive equilibrium of
Definition~\ref{def:mkt:equilibrium} reproduces that allocation, and the equilibrium prices are
\begin{align}
P^R &= p^N,
&
B &= B,
&
P^T &= \alpha B,
&
P^K &= \left(1-\Delta\omega^Y\right)F_K,
\notag\\
Q &= q,
&
P^S &= p^S,
&
P^X &= p^X,
&
s^W &= p^W,
\label{eq:mkt:A:prices}
\end{align}
with $B$ and $p^N$ as in~\eqref{eq:sp:A:Bvalue}, and the market interest rate
equals the planner's, $r=\rho-\dot\lambda/\lambda$.
\end{proposition}

\begin{proof}
Term by term against Section~\ref{sec:sp:A:foc}. We record the six steps because each is a usable
consistency check on a numerical solution.

\emph{(1) Materials.} The final-good firm's condition~\eqref{eq:mkt:setup:foc:goodsfirm} with
$\tau^Y=\Delta\omega^Y$ and $\tau^R=p^W-\zeta$ gives
$P^R=\left(1-\Delta\omega^Y\right)F_R-p^W+\zeta$, which by the planner's production
margin~\eqref{eq:sp:A:FR} is $P^R=p^N$. Hence, by~\eqref{eq:mkt:setup:Bvalue} and
$s^R=d^Wp^P$, $B=p^N+d^Wp^P$, which is the materials
arbitrage~\eqref{eq:sp:A:arbitrage}, so $B=B$.

\emph{(2) Recycling.} Given $B=B$, the recycler's
conditions~\eqref{eq:mkt:setup:foc:recycler} read $P^T=\alpha B$ and
$P^K=a'\!\left(x\right)B$. The latter coincides with the final-good firm's rental
condition $P^K=\left(1-\Delta\omega^Y\right)F_K$ exactly when the planner's capital-allocation
condition~\eqref{eq:sp:A:foc:KR} holds, so the capital market clears at the planner's $x$ and $K^R$.

\emph{(3) Treatment.} Substituting $P^T=\alpha B$, $\tau^E=p^P$ and
$\tau^W=\Delta\omega^W$ into the collector's condition~\eqref{eq:mkt:setup:foc:varpi} gives
$\alpha B+\left(1-d^W\right)p^P=\left(1+\Delta\omega^W\right)dc^T/d\varpi$, which
is~\eqref{eq:sp:A:foc:varpi}. Substituting the same into the zero-profit
condition~\eqref{eq:mkt:setup:zeroprofit} gives
$s^W=\left(1+\Delta\omega^W\right)c^W+\left(1-\varpi+d^W\varpi\right)p^P-\varpi\alpha B$,
which is the planner's waste condition~\eqref{eq:sp:A:foc:W}: $s^W=p^W$.

\emph{(4) Extraction and exploration.} With $\tau^N=\Delta\omega^N$ and
$\tau^{N,S}=\Omega^{N,S}p^W$, the mine's condition~\eqref{eq:mkt:setup:foc:mine} reads
$P^R=\left(1+\Delta\omega^N\right)C^N_N+P^S+\Omega^{N,S}p^W$, and comparison with
$P^R=p^N$ from step~(1) and the definition of $p^N$
in~\eqref{eq:sp:A:Bvalue} gives $P^S=p^S$. The exploration condition then gives
$P^X=p^X$ against~\eqref{eq:sp:A:foc:D}. The mine's costate
equations~\eqref{eq:mkt:setup:costate:mine} coincide with~\eqref{eq:sp:A:costate:S}
and~\eqref{eq:sp:A:costate:X} once $r$ is shown to equal the planner's rate, which is step~(6).

\emph{(5) Consumption and capital.} With $\tau^C=\Delta\omega^C$ the household's
condition~\eqref{eq:mkt:setup:household-euler} is~\eqref{eq:sp:A:foc:C}, which identifies the
household's $\lambda$ with the planner's. With $\tau^K$ as in~\eqref{eq:mkt:A:tauK},
$Q=1+\tau^K=1-\Phi^I\left(1-\chi\right)\left(p^W-p^M\right)=q$ by~\eqref{eq:sp:A:foc:I}.

\emph{(6) The interest rate.} Substituting $P^K$ from step~(2) and $Q=q$
into~\eqref{eq:mkt:setup:household-euler} gives
$r=\left[\left(1-\Delta\omega^Y\right)F_K+\dot q\right]/q-\delta$, which is the planner's capital
costate~\eqref{eq:sp:A:costate:K}. Since both $r$'s are defined as $\rho-\dot\lambda/\lambda$ and the
two $\lambda$'s were identified in step~(5), the discount rates coincide and step~(4) is complete.

\emph{(7) Pollution and the remaining relations.} The household and the firms take $P$ as given, so
$p^P$ has no private counterpart; it enters only through $\tau^E$. Its
law of motion~\eqref{eq:sp:A:costate:P} is the government's, not an equilibrium condition, and the
private output loss from pollution is $\left(1-\tau^Y\right)F_P=\left(1-\Delta\omega^Y\right)F_P$,
which is the term appearing in the planner's damage $d$
in~\eqref{eq:sp:A:damage}. The goods constraint holds by Proposition~\ref{prop:mkt:walras}, the
ledger and the intensity definition hold by Definition~\ref{def:mkt:equilibrium}, and the state
equations are common to the two economies. The transversality
conditions~\eqref{eq:sp:A:tvc} are imposed as the household's and the mine's no-Ponzi conditions.
\end{proof}

The proposition asserts that this instrument set implements \emph{an} optimal allocation. That it
implements \emph{the} optimal allocation requires uniqueness, which Remark~\ref{rem:sp:A:soc} does not
establish; see~\ref{op:sp:soc}.


\subsection{Reading the instrument set}\label{sec:mkt:A:reading}

\smalltitle{The materials block is a deposit--refund scheme.} The two instruments that carry $p^W$ are
a charge of $\tau^R=p^W-\zeta$ per tonne of material entering production and a payment of $s^W=p^W$ per
tonne of waste handed to the collection system. Ignore the wedge $\zeta$ for a moment --- it is zero
under Version~B and small whenever $\chi$ is --- and the pair is a deposit--refund scheme at rate
$p^W$: mass pays on the way in and is refunded on the way out. That the refund equals the deposit is
not an assumption. It is the content of the planner's waste condition~\eqref{eq:sp:A:foc:W}, which
step~(3) of the proof shows is nothing other than the zero-profit condition of a competitive collection
sector facing an emission tax and selling treated waste at $\alpha B$. The shadow price of
waste, in other words, \emph{is} the break-even collection payment.

Two things follow. First, the collection sector is financed, not taxed: $s^W>0$ whenever $p^W>0$, and
the payment is exactly what the sector needs to discharge its obligation given the emission tax it
faces and the revenue it earns from selling treated waste. Second, the deposit is source-neutral. Since
$R=N+R^R$, it can equivalently be collected at the mine and at the recycler rather than at the
final-good firm's gate, and levying it at a different rate on the two sources would break the materials
arbitrage~\eqref{eq:sp:A:arbitrage} directly. Recycled material must bear the same charge per tonne as
virgin material, because recycling does not remove a tonne of waste from the economy --- it defers it
by one circuit, as Section~\ref{sec:sp:A:reading} established. The often-proposed alternative --- a
recycling subsidy financed by a tax on virgin material only --- corresponds to a source-differentiated
deposit and is not optimal here at any rate.

\smalltitle{Why recycling is nonetheless favoured.} The recycler receives $s^R=d^Wp^P$ and the mine
pays $\tau^{N,S}=\Omega^{N,S}p^W$, so the two sources are treated asymmetrically after all --- but for
reasons that have nothing to do with the deposit. Recovery avoids the emission $d^W$ that the tonne
would have caused had it been treated and not recycled, and extraction drags up overburden that
recovery does not. Both are physical asymmetries between the two sources, and both are priced
separately from the mass itself. This is the instrument-level statement of the arbitrage: the
comparison of sources is settled by extraction cost, rent, overburden and avoided emission, and the
waste price cancels from it.

\smalltitle{The capital instrument is the interest on the deposit.} Under
Version~A, \eqref{eq:mkt:A:tauK} makes $\tau^K$ negative whenever $p^W>p^M$, so gross capital formation
is \emph{subsidised}. The reason is now transparent. The deposit is levied at entry, on the
presumption that the tonne is shed at once. A tonne embodied in new capital is not shed at once: it is
shed as the asset depreciates, at a present value of $p^M$ rather than $p^W$ per tonne. The deposit
therefore overcharges stored mass by $p^W-p^M$ per tonne, and $\tau^K$ rebates exactly that, on the
$\Phi^I$ tonnes per unit of gross formation, damped by $\left(1-\chi\right)$ because raising $G$ also
dilutes $\Omega$ and hence lowers $\Phi^I$ itself. It is levied on gross formation because replacement
investment embodies mass on exactly the same terms as net accumulation.

Two consequences are worth recording. The instrument is not Pigouvian in the usual sense: it corrects
no externality of its own, but the timing of a charge levied elsewhere. And its sign is tied to two
others: by~\eqref{eq:sp:A:costate:M} the sign of $p^M$ follows that of $p^W$, so when $p^W$ turns
negative --- the late-transition case of~\eqref{eq:sp:A:pW-sign}, in which waste becomes an asset ---
the deposit becomes a subsidy on material entering production, the refund becomes a charge for handing
waste over, and $\tau^K$ becomes a tax. All three flip together, and a numerical solution in which they
do not is wrong.

\smalltitle{The composition block.} The five instruments $\tau^j=\Delta\omega^j$ have no counterpart
in Version~B and are the market face of the wedge derived in Section~\ref{sec:sp:A:foc}. Their reading
is~(iii) of Section~\ref{sec:mkt:A:what}: an activity that uses the final good dilutes the intensity at
which capital goods embody mass, and so shifts mass from the durable stock into the current waste flow.
The rate is the same for all five up to the relative intensity $\omega^j$, they carry a plus sign on
spending and a plus sign on \emph{output} --- output adds $\omega^YY$ to the waste base, so producing a
unit is worth less than its marginal product --- and all five vanish with $\chi$ or with $p^W-p^M$.

It is worth being explicit that $\tau^C>0$ here is a consumption tax levied on environmental grounds,
which Section~\ref{sec:sp:B:nowedge} argues is generally an artefact. It is not one here, and the test
is the rate: an artefactual wedge would be proportional to $p^W$, this one is proportional to
$\chi\left(p^W-p^M\right)$. The distinction is not academic --- it is the difference between a
consumption tax that survives when postponing waste is worthless and one that does not.


\subsection{The composition block as a single fee on waste}\label{sec:mkt:A:gate}

The five composition instruments look unwieldy: ad valorem taxes on output, on consumption, on
extraction effort, on exploration effort, and on the waste sector's own expenditure, each at its own
rate. They are in fact one instrument, and the natural one.

\begin{proposition}[The composition block is a gate fee on waste generation]\label{prop:mkt:A:gate}
Suppose the collection sector charges a fee $g$ per tonne to whoever hands waste over, on the base
$\Omega\mathcal D$ --- the waste attributable to final-good spending --- and $s^W$ adjusts through
zero profit. Then setting
\begin{align}
g &= \zeta = \chi\left(p^W-p^M\right),
&
\tau^j &= 0\ \ \text{for all } j\in\left\{Y,C,N,D,W\right\},
&
s^W &= p^W-\zeta\,\frac{\Omega\mathcal D}{W},
\label{eq:mkt:A:gate}
\end{align}
and leaving the remaining instruments at their values in
Proposition~\ref{prop:mkt:A:implementation}, implements the same allocation.
\end{proposition}

\begin{proof}
Spending one unit of the final good on activity $j$ generates $\Omega\omega^j$ tonnes of the fee base,
so the effective cost of that unit becomes $1+g\Omega\omega^j$, and producing one unit of output
generates $\Omega\omega^Y$ tonnes, so its effective value becomes $1-g\Omega\omega^Y$. With $g=\zeta$
these are $1+\Delta\omega^j$ and $1-\Delta\omega^Y$, by $\Delta=\zeta\Omega$
in~\eqref{eq:mkt:A:zeta}, which are the factors the ad valorem instruments produced. Every first-order
condition in the proof of Proposition~\ref{prop:mkt:A:implementation} is therefore unchanged. The
collection sector's zero-profit condition~\eqref{eq:mkt:setup:zeroprofit} acquires the fee revenue
$g\Omega\mathcal D$ on the left and the fee it pays on its own expenditure on the right, which is the
$\left(1+\Delta\omega^W\right)$ factor already there; solving for $s^W$ gives the stated value.
\end{proof}

The proposition matters for three reasons.

First, it says that the instrument the coarse intuition recommends --- make each activity pay for the
waste it generates --- is correct \emph{in form} under Version~A. What the coarse intuition gets wrong
is the rate. The fee is not $p^W$, the social cost of a tonne of waste; it is $\zeta$, the value of the
storage margin, which is smaller by the factor $\chi$ and smaller again by the difference between
shedding a tonne now and shedding it later. A fee set at $p^W$ would overcharge every activity in the
economy by $\Omega\omega^j\left(p^W-\zeta\right)$ per unit, which is precisely the spurious wedge
of~\eqref{eq:sp:B:spurious-C}.

Second, it makes the deposit--refund structure exact. Under~\eqref{eq:mkt:A:gate} every tonne pays
$\zeta$ at the gate and $p^W-\zeta$ at entry, and is refunded $p^W-\zeta$ on collection --- so the
deposit and the refund are equal, at the rate $p^W-\zeta$, and the only genuine \emph{disposal} charge
in the optimal system is $\zeta$. The wedge between deposit and refund noted in
Section~\ref{sec:mkt:A:reading} was the composition block in disguise.

Third, the base matters and is not the obvious one. The fee must fall on the waste generated by
final-good spending and \emph{not} on the mass released by depreciation, $\delta M^K$, nor on the
displaced mass $\mathcal N$. The reason is that the fee is a dilution charge, not a waste charge: only
final-good spending enters $\mathcal D$ and moves $\Omega$. Demolition debris pays no fee because
scrapping a building dilutes nothing; it merely discharges a liability whose deposit was paid when the
mass entered, and which $\tau^K$ has already priced. This is the market counterpart of the observation
after~\eqref{eq:sp:A:costate:M} that the embodied-material costate is the one equation in the system
carrying no wedge. Extending the fee to demolition debris would require the household to hold $M^K$ as
a private state with its own liability price, and would have to be undone in $\tau^K$.


\subsection{The complete system, and what the government must know}\label{sec:mkt:A:system}

\smalltitle{The system.} A competitive equilibrium under Version~A consists of: the five state
equations~\eqref{eq:sp:A:states-motion}; the household's two
conditions~\eqref{eq:mkt:setup:household-euler}; the final-good firm's
two~\eqref{eq:mkt:setup:foc:goodsfirm}; the recycler's two~\eqref{eq:mkt:setup:foc:recycler}; the
mine's two static and two dynamic conditions~\eqref{eq:mkt:setup:foc:mine}--\eqref{eq:mkt:setup:costate:mine};
the collector's~\eqref{eq:mkt:setup:foc:varpi} and~\eqref{eq:mkt:setup:zeroprofit}; market
clearing~\eqref{eq:mkt:setup:clearing}; the goods constraint; the ledger and the intensity definition,
solved jointly for $W$ and $\Omega$ as in Section~\ref{sec:sp:A:system}; and the government
budget~\eqref{eq:mkt:setup:govbudget}. Given instrument paths this system determines the equilibrium
for \emph{any} policy, optimal or not, which is worth stating separately: it is the object a
counterfactual exercise needs, and Part~\ref{part:sp} does not contain it. What it does not deliver is
the welfare ranking of two such policies, which still requires evaluating $U_0$; see~\ref{op:sp:nonopt}.

\smalltitle{Informational requirements.} To set the instruments of
Proposition~\ref{prop:mkt:A:implementation} the government must know $p^W$, $p^P$, $p^M$, $\Omega$,
$\chi$ and every relative intensity $\omega^j$. The last two are the demanding ones. $\Omega$ and the
$\omega^j$ are objects of the material-flow accounts, measured with considerable error, and under
Version~A they enter every instrument in the lower block of Table~\ref{tab:mkt:instruments} and, through
$\Delta$, the implicit determination~\eqref{eq:sp:A:foc:W-solved} of $p^W$ itself. There is no part of
the optimal policy that can be computed without them. Section~\ref{sec:mkt:B:compare} contrasts this
with Version~B, where the entire instrument set is free of them.

\begin{remark}[Integrating the recycler]\label{rem:mkt:A:integration}
Because the recycling sector earns exactly zero profit at every date, it can be merged into the
final-good firm --- which then buys treated waste at $P^T$, allocates capital between production and
recovery, and uses $N+\mathcal R\left(T,K^R\right)$ as its material input --- without changing any
equilibrium condition or any instrument. The merged firm's internal shadow price of secondary material
is $B$. The separation is kept because it makes $P^T$ and $B$ observable objects rather than internal
ones, and because the deposit $\tau^R$ then has a visible base.
\end{remark}

\begin{remark}[Where $\zeta$ acts]\label{rem:mkt:A:zeta}
The multiplier $\zeta$ enters the instrument set at exactly two points: it reduces the deposit
$\tau^R$ below $p^W$, and it scales the composition block through $\Delta=\zeta\Omega$. Under Version~B
the mass-constraint multiplier $\eta$ takes over the first role and there is no second one; this is the
sense in which Section~\ref{sec:sp:B:compare}'s statement that ``$\eta$ is replaced by $\zeta$''
operates in the market economy as well.
\end{remark}
